%%% Основные сведения %%%
\newcommand{\thesisAuthorLastName}{\fixme{Макаров}}
\newcommand{\thesisAuthorOtherNames}{\fixme{Станислав Олегович}}
\newcommand{\thesisAuthorInitials}{\fixme{С.\,О.}}
\newcommand{\thesisAuthor}             % Диссертация, ФИО автора
{%
    \texorpdfstring{% \texorpdfstring takes two arguments and uses the first for (La)TeX and the second for pdf
        \thesisAuthorLastName~\thesisAuthorOtherNames% так будет отображаться на титульном листе или в тексте, где будет использоваться переменная
    }{%
        \thesisAuthorLastName, \thesisAuthorOtherNames% эта запись для свойств pdf-файла. В таком виде, если pdf будет обработан программами для сбора библиографических сведений, будет правильно представлена фамилия.
    }
}
\newcommand{\thesisAuthorShort}        % Диссертация, ФИО автора инициалами
{\thesisAuthorInitials~\thesisAuthorLastName}
%\newcommand{\thesisUdk}                % Диссертация, УДК
%{\fixme{xxx.xxx}}
\newcommand{\thesisTitle}              % Диссертация, название
%{\fixme{Совершенствование методики передачи координат с использованием метода PPP}}
{\fixme{Совершенствование методики применения метода PPP при построении геодезических сетей}}
\newcommand{\thesisSpecialtyNumber}    % Диссертация, специальность, номер
{\fixme{1.6.22}}
\newcommand{\thesisSpecialtyTitle}     % Диссертация, специальность, название (название взято с сайта ВАК для примера)
{\fixme{Геодезия}}
%% \newcommand{\thesisSpecialtyTwoNumber} % Диссертация, вторая специальность, номер
%% {\fixme{XX.XX.XX}}
%% \newcommand{\thesisSpecialtyTwoTitle}  % Диссертация, вторая специальность, название
%% {\fixme{Теория и~методика физического воспитания, спортивной тренировки,
%% оздоровительной и~адаптивной физической культуры}}
\newcommand{\thesisDegree}             % Диссертация, ученая степень
{\fixme{кандидата технических наук}}
\newcommand{\thesisDegreeShort}        % Диссертация, ученая степень, краткая запись
{\fixme{канд.~техн.~наук}}
\newcommand{\thesisCity}               % Диссертация, город написания диссертации
{\fixme{Новосибирск}}
%
\newcommand{\nextyear}{\the\numexpr\the\year+1\relax}  % следующий год
%
\newcommand{\thesisYear}               % Диссертация, год написания диссертации
{\the\year}							   % текующий год
%{\nextyear}			  				% следующий год ;)
\newcommand{\thesisOrganization}       % Диссертация, организация
{\fixme{МИНИСТЕРСТВО СЕЛЬСКОГО ХОЗЯЙСТВА РОССИЙСКОЙ ФЕДЕРАЦИИ\\Федеральное государственное бюджетное образовательное\\учреждение высшего образования\\<<ГОСУДАРСТВЕННЫЙ УНИВЕРСИТЕТ ПО ЗЕМЛЕУСТРОЙСТВУ>>\\ГУЗ}}
\newcommand{\thesisOrganizationShort}  % Диссертация, краткое название организации для доклада
{\fixme{ГУЗ}}

\newcommand{\thesisInOrganization}     % Диссертация, организация в предложном падеже: Работа выполнена в ...
{\fixme{Федеральном государственном бюджетном образовательном учреждении высшего образования «Государственный университет по землеустройству»}}

%\newcommand{\supervisorDead}{}           % Рисовать рамку вокруг фамилии
\newcommand{\supervisorFio}              % Научный руководитель, ФИО
{\fixme{Тихонов Александр Дмитриевич}}
\newcommand{\supervisorRegalia}          % Научный руководитель, регалии
{\fixme{кандидат технических наук, доцент }}		 % 
\newcommand{\supervisorFioShort}         % Научный руководитель, ФИО
{\fixme{А.\,Д.~Тихонов}}
\newcommand{\supervisorRegaliaShort}     % Научный руководитель, регалии
{\fixme{к.т.н.}}

%% \newcommand{\supervisorTwoDead}{}        % Рисовать рамку вокруг фамилии
%% \newcommand{\supervisorTwoFio}           % Второй научный руководитель, ФИО
%% {\fixme{Фамилия Имя Отчество}}
%% \newcommand{\supervisorTwoRegalia}       % Второй научный руководитель, регалии
%% {\fixme{уч. степень, уч. звание}}
%% \newcommand{\supervisorTwoFioShort}      % Второй научный руководитель, ФИО
%% {\fixme{И.\,О.~Фамилия}}
%% \newcommand{\supervisorTwoRegaliaShort}  % Второй научный руководитель, регалии
%% {\fixme{уч.~ст.,~уч.~зв.}}

\newcommand{\opponentOneFio}           % Оппонент 1, ФИО
{\fixme{Брынь Михаил Ярославович}}
\newcommand{\opponentOneRegalia}       % Оппонент 1, регалии
{\fixme{доктор технических наук, профессор}}
\newcommand{\opponentOneJobPlace}      % Оппонент 1, место работы
{\fixme{Федеральное государственное бюджетное учреждение высшего образование <<Санкт-Петербургский государственный университет путей сообщений имени императора Александра I>>}}
\newcommand{\opponentOneJobPost}       % Оппонент 1, должность
{\fixme{заведующий кафедрой инженерной геодезии}}

\newcommand{\opponentTwoFio}           % Оппонент 2, ФИО
{\fixme{Косарев Николай Сергеевич}}
\newcommand{\opponentTwoRegalia}       % Оппонент 2, регалии
{\fixme{кандидат технических наук}}
\newcommand{\opponentTwoJobPlace}      % Оппонент 2, место работы
{\fixme{Федеральное государственное бюджетное образовательное учреждение высшего образования <<Сибирский государственный университет геосистем и технологий>>}}
\newcommand{\opponentTwoJobPost}       % Оппонент 2, должность
{\fixme{доцент кафедры инженерной геодезии и маркшейдерского дела}}

%% \newcommand{\opponentThreeFio}         % Оппонент 3, ФИО
%% {\fixme{Фамилия Имя Отчество}}
%% \newcommand{\opponentThreeRegalia}     % Оппонент 3, регалии
%% {\fixme{кандидат физико-математических наук}}
%% \newcommand{\opponentThreeJobPlace}    % Оппонент 3, место работы
%% {\fixme{Основное место работы c длинным длинным длинным длинным названием}}
%% \newcommand{\opponentThreeJobPost}     % Оппонент 3, должность
%% {\fixme{старший научный сотрудник}}

\newcommand{\leadingOrganizationTitle} % Ведущая организация, дополнительные строки. Удалить, чтобы не отображать в автореферате
{~} %%{\fixme{Федеральное государственное бюджетное образовательное учреждение высшего образования «ПГУПС императора
	%%	Александра I» (предварительно)
%%}}

\newcommand{\defenseDate}              % Защита, дата
{\fixme{DD MMMM 2026~г.~в~12-00 часов}}
\newcommand{\defenseCouncilNumber}     % Защита, номер диссертационного совета
{\fixme{24.2.402.01}}
\newcommand{\defenseCouncilTitle}      % Защита, учреждение диссертационного совета
{\fixme{ФГБОУ ВО <<Сибирский государственный университет геосистем и технологий>>}}
\newcommand{\defenseCouncilAddress}    % Защита, адрес учреждение диссертационного совета
{\fixme{630108, г.~Новосибирск, ул.~Плахотного, 10, ауд.~402}}
\newcommand{\defenseCouncilPhone}      % Телефон для справок
{\fixme{+7~(383)~343-39-23}}

\newcommand{\defenseSecretaryFio}      % Секретарь диссертационного совета, ФИО
{\fixme{Аврунев Евгений Ильич}}
\newcommand{\defenseSecretaryRegalia}  % Секретарь диссертационного совета, регалии
{\fixme{канд.~техн.~наук, доц.}} % Для сокращений есть ГОСТы, например: ГОСТ Р 7.0.12-2011 + http://base.garant.ru/179724/#block_30000

\newcommand{\thesisHref} % Ссылка на текст диссертации в интернете
{https://sgugit.ru/science-and-innovations/dissertation-councils/makarov-stanislav-olegovich/}

\newcommand{\synopsisLibrary}          % Автореферат, название библиотеки
{\fixme{и на сайте \defenseCouncilTitle \\\href{\thesisHref}{\thesisHref} }}

\newcommand{\synopsisDate}             % Автореферат, дата рассылки
{\fixme{DD mmmmmmmm} 
\the\year % текущий год
%\nextyear % следующий год
~года}

% To avoid conflict with beamer class use \providecommand
\providecommand{\keywords}%            % Ключевые слова для метаданных PDF диссертации и автореферата
{метод PPP, геодезические сети, высокоточное позиционирование, ГНСС, передача координат, точность определения, эфемериды, PPP, тропосферная задержка, ионосферная задержка, Российская Федерация,
GNSS; Precise Point Positioning; точность позиционирования; спутниковые методы; PDOP-фактор;
государственная геодезическая сеть
}
% Восстанавливаем макрос \No, который был удалён из Бабеля в апреле 2013.
%\newcommand{\No}{\textnumero}
% Глобально боремся с оверфулами, но нужно следить за переразреженнеми строками...
%\sloppy
% Запрет переносов в словах, начинающихся с Прописной буквы
%\uchyph=0

 
\newcommand{\colophon}    % Выходные данные автореферата -- последняя страница и оборот титула
{\fixme{
	Изд. лиц. ЛР № 020461 от 04.03.1997. 
\\
Подписано в печать \blank[\widthof{999}].\blank[\widthof{999}].\blank[\widthof{99999}]. 
Формат 60\(\times\)84 1/16. 
\\
Печ. л. 1,0. 
Тираж 100 экз.
Заказ № \blank[\widthof{999999999999}]
\\
Редакционно-издательский отдел СГУГиТ
\\
630108, Новосибирск, Плахотного, 10. 
\\
Отпечатано в картопечатной лаборатории СГУГиТ
\\
630108, Новосибирск, Плахотного, 8.
}}